% =====================================================================
%  II. Background and Related Work
%
%  Ba viec muc nay phai lam, theo dung thu tu uu tien cua reviewer:
%    R1-1  : them QMI 2026 (bai R1 chi dich danh) + cap nhat Table I
%    R3-5  : dinh vi ro so voi arXiv:2403.07059 va 2409.04406
%    tu giac: trich Gillani et al. (2608.18155) -- bai dong thoi, trung
%             ca hai dataset. Nop sau ho 2 thang ma lo di la reviewer tu tim ra.
%
%  Giong van: KHONG cai nhau voi R3. Neu ho noi benchmark khong phai dong
%  gop, ta neu su that la hai bai chinh ho tin cay cung la benchmark thuan
%  tuy -- roi di tiep.
% =====================================================================

\subsection{Quantum kernels}
\revnote{rewrite}{Submitted version had Definitions 1--2 and Propositions 1--2; merged, and Proposition 1 demoted to a cited fact (F1) as R3-3 asks.}
\begin{revbody}

A quantum kernel embeds $\bm{x}$ into a state $|\phi(\bm{x})\rangle$ prepared by
a parameterised circuit and reads off the fidelity as a similarity
\cite{havlicek2019,schuld2019}. The construction sits inside a broader quantum
machine learning programme \cite{ref5,ref6,ref7}, and the identification of
parameterised-circuit models with kernel machines \cite{ref9} is what makes the
quantum component substitutable into a classical learner. Separations have been
proved for constructed problems \cite{ref18}, but the conditions under which they
extend to data that arise in practice are restrictive \cite{ref19,ref20}.
Two results bound expectations. Alignment-based
bounds \cite{cristianini2002,cortes2012} say that a kernel better aligned with
the labels admits a tighter generalisation guarantee at equal empirical risk,
but say nothing about whether a given circuit produces such a kernel. In the
other direction, fidelity kernels from expressive maps concentrate exponentially
toward a constant as the qubit count grows \cite{thanasilp2024}, an effect
distinct from, though related to, the barren plateaus of variational training
\cite{ref10}, so the same expressivity \cite{ref21} that makes a map interesting
makes its Gram matrix uninformative at width. Sec.~\ref{sec:res-width} measures both effects here: the alignment gain
is large, and the concentration is what eventually destroys it.
\end{revbody}

\subsection{Quantum machine learning for intrusion detection}
\revnote{rewrite}{Adds the 2026 survey and Cirillo et al. (QMI 2026), the study Reviewer 1 named in R1-1.}
\begin{revbody}

Applications of quantum kernels and variational classifiers to network intrusion
detection have grown quickly \cite{ref13,ref14,ref15}, most using NSL-KDD
\cite{ref4} or a successor corpus \cite{ref25}; a 2026 survey \cite{kaissar2026} covers the field through
January of that year and reaches two conclusions that bear directly on this
paper: most published approaches rely on simulated quantum environments and
legacy datasets, and evaluation practice is inconsistent across studies. The
recurring pattern is a single operating point (one qubit count, one
training-set size, one class prior) with an accuracy compared against a small
set of classical baselines.

The most directly comparable study is the benchmark of Cirillo \emph{et al.}
\cite{qmi2026}, which evaluates Pegasos-QSVC, a variational classifier and a
hybrid quantum--classical network on ToN\_IoT and NSL-KDD under ideal and three
backend-derived noise models, and reports Pegasos-QSVC best at $94.60\%$
accuracy and $94.13\%$ $F_1$ on NSL-KDD.

Two of their findings bear on ours. They compare against SVM, XGBoost, random
forest and a neural network at matched dimensionality and conclude that the
classical models retain higher peak accuracy, XGBoost reaching $99.1\%$ on
ToN\_IoT, while the quantum models stay competitive in low-dimensional
feature spaces; our own ranking agrees (Sec.~\ref{sec:res-c2}). And they sweep
the qubit count from $2$ to $10$ against two-qubit gate count, finding
performance concentrated in a $20$--$35$ CNOT band and degrading beyond about
$60$. Our operating point of $24$ CNOTs sits inside that band, and the collapse
we report at $112$ is consistent with their threshold.

The two studies diverge on the mechanism, and that divergence is part of our
contribution: they attribute the degradation to two-qubit gate error
accumulating on noisy backends, whereas we observe the same collapse under
\emph{exact, noiseless} simulation and trace it to concentration of the Gram
matrix. Hardware noise is not a necessary ingredient; the bound is intrinsic to
the kernel.

What that study does not vary is the training-set size, fixed at a single $15\%$
class-balanced subsample ($15{,}114$ training instances for NSL-KDD) drawn under
a random $80/20$ split rather than the official partitions.
\end{revbody}

\subsection{Benchmarking studies of quantum models}
\revnote{new}{New subsection. The two benchmarking studies Reviewer 3 cited as prior evidence (R3-5), read in full.}
\begin{revbody}

Two broad benchmarks are relevant precisely because they are sceptical. Bowles,
Ahmed and Schuld \cite{bowles2024} evaluate twelve QML models over $160$ datasets
(predominantly synthetic, so that data properties can be controlled) in
noise-free statevector simulation at up to $18$ qubits, against an MLP, an
RBF-kernel SVM and a CNN, and find that quantum models rarely beat well-tuned
classical baselines. Schnabel and Roth \cite{schnabel2024} run over $20{,}000$
tuned models comparing fidelity and projected quantum kernels across five
dataset families, at up to $15$ qubits against RBF-kernel baselines, and use
correlation analysis to identify which design choices matter. Both conclude that
general-purpose feature maps seldom win once the classical comparison is done
properly, and our results agree with them at the aggregate level.

Neither varies the training-set size: Bowles \emph{et al.} subsample fixed
counts per task, and Schnabel and Roth fix $M=240$ training and $M'=60$ test
points throughout.

Concurrently with this revision, Gillani \emph{et al.} \cite{gillani2026}
benchmark quantum and classical models for intrusion detection across four
corpora, including both of ours, at eight qubits with a $4/6/8/12$-qubit
ablation, against five tuned classical baselines including a random forest and
XGBoost, with Benjamini--Hochberg control over $108$ comparisons on NSL-KDD.
Their conclusion, that reported quantum gains are largely attributable to the
classical preprocessing front-end once the same front-end is given to classical
models, stands in a relation to Sec.~\ref{sec:res-c4} that is closer to a
discrepancy than to corroboration, and we report it as such. Our positive
result is obtained under exactly the matched front-end they describe: every
classical baseline consumes the same four components as the circuit, and the
kernel still leads the tree ensembles at $N\geq5000$. Two differences may
account for it, and the data cannot separate them. Their comparison is made at
one training-set size, the full $125{,}973$-record partition, where our
reference arm shows full-feature XGBoost at $0.804$ above any four-qubit model;
ours sweeps $N$ and finds the lead only between $5000$ and $10^{4}$ records, on
an embedding of four dimensions. And their front-end is given to classical
models at eight qubits' worth of features, where Sec.~\ref{sec:res-width}
shows our kernel has already collapsed. Read together, the two studies agree
that the front-end carries most of what is reported as quantum advantage, and
disagree on whether a residual lead exists in a narrow window of training-set
size and embedding width that their design does not visit. We do not claim
breadth over that study: they evaluate four datasets to our two, at eight
qubits to our four and six, with a stricter error-rate criterion.
\end{revbody}

\subsection{Bounded advantage windows in adjacent domains}
\revnote{new}{New subsection. Carducci, ICAD 2026, the study Reviewer 4 named in R4-1.}
\begin{revbody}

The question this paper asks is being asked elsewhere, and in at least one
adjacent domain it has been given the same shape. Carducci \cite{carducci2026}
opens by arguing that malware analysts ``do not need another general claim that
quantum computing is either broadly superior or broadly impractical'' but rather
``need to know which kinds of malware representations justify quantum effort''
which is, on a different task, the reframing this revision makes.

That study places malware representations on an eleven-rank
representational-complexity scale and reports that the measured advantage is
confined to an \emph{intermediate} window: at the co-occurrence and local-order
representations the quantum workflow retrieves malware where the classical
baseline retrieves none ($4.29\%$ against $0\%$), while classical methods remain
stronger both at the lower-complexity statistical representations and at the
higher-complexity ones. The advantage survives parameter tuning, which the
author reports specifically to rule out a favourable-setting artefact, and the
stated practical implication is to prefer classical methods wherever mature
inductive biases already exist.

Three features of that result recur in ours, on a different task and with a
different quantum primitive: the advantage is bounded on \emph{both} sides
rather than growing with capability; the authors test explicitly whether
re-tuning explains it, as our frozen and tuned arms do; and the region where the
quantum method helps is characterised by \emph{local pairwise structure}, which
is what the $\binom{n}{2}$ pairwise phase terms of the \ZZ{} map encode and what
Sec.~\ref{sec:res-width} shows to stop paying off once $n$ grows.

We could not obtain that full text through our institutional access, so the
comparison rests on its abstract, introduction and index terms; we could not
establish its qubit count or whether it varies training-set size. We present the
parallel as suggestive rather than as evidence, and do not tabulate the study in
Table~\ref{tab:novelty}, whose cells are read from full texts.
\end{revbody}

\subsection{Positioning}
\revnote{rewrite}{Table I rebuilt: from 4 studies and 3 columns to 5 studies and 13 rows, every cell read from the full text (AE-2, R3-5).}
\input{tables/novelty_matrix}
\begin{revbody}


Table~\ref{tab:novelty} places this work against those four studies on the axes
that distinguish them; each cell was read from the full text rather than from an
abstract. The distinguishing axis is not the model, the dataset, the circuit
width or the noise treatment, on all of which we are comparable or narrower;
two of the four vary the qubit count, and two compare against tree ensembles. It
is that \emph{none of the four sweeps the training-set size}, which each fixes at
a single value, and the training-set size is the variable that decides the answer
here: the ordering between \QSVM{} and every strong classical baseline reverses
inside the range $N\in[10^{2},10^{4}]$ under the natural class prior. A benchmark
conducted at a single $N$, as all four are and as our own submitted version
was, can therefore report either outcome and be internally consistent.

We take the reviewers' point that benchmarking alone is not a contribution, and
we note that the two studies most often cited as establishing what is already
known \cite{bowles2024,schnabel2024} are themselves benchmarks without a new
kernel, feature map, algorithm, hardware demonstration or theoretical result.
What makes a benchmark worth publishing is not its novelty of apparatus but
whether it changes what the field believes. We report an ordering that reverses
with sample size, a selection rule whose output transfers across datasets, and a
measured mechanism, Gram-matrix concentration, that predicts where the
method stops working. The last of these is the closest we come to a theoretical
contribution, and we present it as an empirical regularity rather than as a
theorem, which is what our evidence supports.
\end{revbody}
